\chapter{Mise en place du modèle d'IA générative}

% ─────────────────────────────────────────────────────────────────────────────
\section{Introduction}

Ce chapitre décrit la phase de préparation des données et de mise en place
des modèles d'intelligence artificielle pour chacun des quatre modules du système.

Pour les modules conversationnels fondés sur une architecture RAG —
Noa et le Support Assistant — cette phase consiste en la construction
et l'enrichissement des bases de connaissances documentaires exploitées
lors de la recherche par similarité sémantique.

Pour le classificateur de harcèlement, elle couvre le cycle complet
d'apprentissage automatique : collecte des données, nettoyage et prétraitement,
augmentation, labeling, entraînement et évaluation comparative des stratégies.

Le module de soutien émotionnel, dont l'architecture repose exclusivement sur
l'ingénierie de prompt sans base documentaire ni données d'entraînement,
n'est pas concerné par cette phase.

% ─────────────────────────────────────────────────────────────────────────────
\section{Préparation des bases de connaissances RAG}

Les modules Noa et Support Assistant reposent sur une architecture
\textit{Retrieval-Augmented Generation} (RAG) dont la qualité des réponses
dépend directement de la qualité et de la structure des documents indexés.
La phase de préparation consiste à transformer les sources documentaires brutes
en fichiers Markdown structurés et enrichis, exploitables par la base
vectorielle ChromaDB.

\subsection{Module Noa — sources documentaires publiques}

\subsubsection{Extraction et conversion}

La base de connaissances de Noa est constituée à partir du contenu du site public
de Netethic.
Les vingt-et-une pages du site ont été extraites et converties au format Markdown.
Cette conversion produit des fichiers textuels structurés,
plus adaptés au découpage en segments que les sources HTML brutes.

\subsubsection{Enrichissement des fichiers Markdown}

Après conversion, chaque fichier Markdown a été relu et enrichi
pour corriger les problèmes courants liés à l'extraction automatique~:

\begin{itemize}
    \item \textbf{Clarification des sous-titres :} les intitulés de section trop courts
          ou ambigus ont été reformulés pour décrire explicitement le contenu
          de la rubrique et améliorer la précision de la recherche par similarité.
    \item \textbf{Complétion des informations manquantes :} certains passages
          incomplets (tarifs, conditions d'accès, coordonnées) ont été complétés
          à partir des sources officielles.
    \item \textbf{Suppression des éléments non informatifs :} menus de navigation,
          pieds de page et formulaires ont été retirés pour ne conserver
          que le contenu textuel utile à la réponse.
\end{itemize}

\subsubsection{Stratégie d'indexation parent-enfant}

Les documents enrichis ont été découpés selon une stratégie parent-enfant~:
les segments courts (enfants) servent à la recherche par similarité sémantique,
tandis que les segments longs (parents) sont transmis au LLM pour générer
une réponse contextuellement riche.
Cette stratégie permet d'allier précision de la recherche et richesse
du contexte fourni au modèle de génération.

Le Tableau~\ref{tab:kb_noa} présente les statistiques de la base de connaissances
constituée pour Noa.

\vspace{6pt}
\begin{table}[H]
\centering
\caption{Statistiques de la base de connaissances — Module Noa}
\label{tab:kb_noa}
\begin{tabular}{|l|c|}
\hline
\textbf{Paramètre} & \textbf{Valeur} \\
\hline
Nombre de pages indexées    & 21 \\
\hline
Nombre de segments (chunks) & 177 \\
\hline
Modèle d'embedding          & \texttt{paraphrase-multilingual-mpnet-base-v2} \\
\hline
Seuil de pertinence         & 0,35 \\
\hline
Base vectorielle            & ChromaDB \\
\hline
\end{tabular}
\end{table}

\subsection{Module Support Assistant — documentation interne}

\subsubsection{Sources documentaires}

La base de connaissances du Support Assistant est constituée à partir
de la documentation technique interne fournie par Smart-Conseil.
Elle couvre les onze catégories fonctionnelles de la plateforme Netethic~:
gestion des utilisateurs, paramétrage, tableaux de bord, rapports, notifications,
intégration, sécurité, conformité RGPD, support technique, administration
et facturation.

\subsubsection{Transformation et structuration}

Les documents fournis au format PDF ont été convertis au format Markdown.
Chaque fichier Markdown a ensuite été structuré par catégorie fonctionnelle,
en suivant le même processus d'enrichissement que pour la base de Noa~:
clarification des sous-titres, complétion des informations manquantes
et suppression des éléments non informatifs.

La collection ChromaDB du Support Assistant est organisée en onze sous-collections,
une par catégorie fonctionnelle.
Cette organisation permet un filtrage précis lors de la recherche vectorielle~:
lorsqu'une catégorie est sélectionnée par l'utilisateur, la recherche
est restreinte aux documents de cette catégorie uniquement,
ce qui améliore la pertinence des résultats retournés.

% ─────────────────────────────────────────────────────────────────────────────
\section{Mise en place du classificateur de harcèlement}

Le classificateur de harcèlement est le module dont la mise en place
diffère fondamentalement des modules conversationnels~:
il ne repose pas sur un pipeline RAG mais sur un cycle complet
d'apprentissage automatique, depuis la collecte des données
jusqu'à l'évaluation comparative des stratégies de classification.

\subsection{Collecte des données}

Le corpus d'expérimentation est constitué de \textbf{977 textes multilingues}
(français et anglais), issus d'une application mobile de modération de contenu.
Tous les textes sont initialement étiquetés \textit{négatif} par l'application.
Le Tableau~\ref{tab:corpus_classif} décrit les caractéristiques du jeu de données brut.

\vspace{6pt}
\begin{table}[H]
\centering
\caption{Caractéristiques du corpus brut — Classificateur de harcèlement}
\label{tab:corpus_classif}
\small
\begin{tabular}{|l|l|}
\hline
\textbf{Caractéristique} & \textbf{Valeur} \\
\hline
Nombre de textes        & 977 \\
\hline
Langues                 & Français et anglais (mélangés) \\
\hline
Source                  & Application mobile de modération de contenu \\
\hline
Encodage                & Latin-1 (ISO-8859-1) \\
\hline
Étiquette initiale      & \textit{negative} (classification de l'application) \\
\hline
\end{tabular}
\end{table}

\subsection{Nettoyage et prétraitement}

L'analyse du corpus a révélé cinq catégories de bruits textuels~:
entités HTML (8 occurrences), images Markdown (3),
textes dépassant 128~caractères (152),
mentions d'utilisateurs (14) et URLs (5).

Un pipeline de prétraitement en cinq étapes a été développé pour normaliser
les textes avant inférence.
Le Tableau~\ref{tab:preprocessing_steps} décrit chaque étape.

\vspace{6pt}
\begin{table}[H]
\centering
\caption{Pipeline de prétraitement du corpus}
\label{tab:preprocessing_steps}
\small
\setlength{\tabcolsep}{4pt}
\begin{tabular}{|c|p{4cm}|p{8cm}|}
\hline
\textbf{Étape} & \textbf{Traitement} & \textbf{Description} \\
\hline
1 & Décodage HTML
  & Conversion des entités HTML (\texttt{\&amp;}, \texttt{\&\#39;}, etc.)
    en caractères lisibles. \\
\hline
2 & Suppression des balises
  & Élimination des balises HTML résiduelles (\texttt{<br>}, \texttt{<a href=...>}, etc.). \\
\hline
3 & Normalisation des mentions
  & Remplacement de toutes les mentions d'utilisateurs par un jeton générique. \\
\hline
4 & Normalisation des URLs
  & Remplacement de toutes les URLs par un jeton générique. \\
\hline
5 & Normalisation des espaces
  & Suppression des sauts de ligne multiples et collapsage des espaces superflus. \\
\hline
\end{tabular}
\end{table}

Le prétraitement conserve l'intégralité des 977~textes (aucune ligne vide
après nettoyage) et ne modifie que 2~prédictions sur 977 (0,2\,\%),
confirmant que les erreurs observées sont intrinsèques au modèle
et non aux données brutes.

\subsection{Augmentation des données}

Le corpus initial présente un déséquilibre de classes~:
après la phase de labeling (voir Section~\ref{subsec:labeling}),
647~textes sont confirmés dans la catégorie négative contre 330~textes
reclassifiés dans des catégories positives ou neutres.
Pour atténuer ce déséquilibre et améliorer la robustesse du modèle,
une stratégie d'augmentation des données a été appliquée à la classe minoritaire.

L'augmentation repose sur la paraphrase automatique par LLM~:
chaque texte de la classe minoritaire est soumis au modèle indépendant
avec une instruction de paraphrase, produisant une version reformulée
sémantiquement équivalente mais lexicalement différente.
Cette technique élargit la diversité des exemples d'entraînement
sans nécessiter de collecte de nouvelles données.

\subsection{Stratégie de labeling — LLM comme juge indépendant}
\label{subsec:labeling}

Pour évaluer la qualité des étiquettes produites par le modèle de classification
initial, une stratégie de \textit{LLM-as-a-Judge} a été adoptée~:
les 977~textes sont soumis à un modèle indépendant (\textbf{Groq API},
modèle \texttt{llama-3.1-8b-instant}) pour obtenir une annotation de référence.

Cette approche évite l'évaluation circulaire dans laquelle un modèle validerait
ses propres prédictions — biais qui rendrait les métriques sans valeur.

Le Tableau~\ref{tab:labeling_results} présente les résultats obtenus.

\vspace{6pt}
\begin{table}[H]
\centering
\caption{Résultats du labeling indépendant — LLM-as-a-Judge}
\label{tab:labeling_results}
\small
\begin{tabular}{|l|c|}
\hline
\textbf{Métrique} & \textbf{Valeur} \\
\hline
Textes confirmés dans la catégorie négative   & 647 / 977 \\
\hline
Textes reclassifiés (positif ou neutre)       & 330 / 977 \\
\hline
Précision du modèle initial                   & 66,2\,\% \\
\hline
Taux d'accord entre les deux modèles         & 63,4\,\% (619 / 977) \\
\hline
Impact du prétraitement sur les prédictions   & 0,2\,\% (2 / 977) \\
\hline
\end{tabular}
\end{table}

Les 33,8\,\% de désaccords s'expliquent principalement par un écart de domaine~:
le modèle initial a été entraîné sur des données formelles anglophones,
alors que le corpus contient de l'argot jeune francophone
(\og trop sous coter \fg{}, \og le meilleurrrrrr \fg{})
que ce modèle interprète incorrectement comme négatif.

\subsection{Entraînement et optimisation}

Quatre stratégies de classification ont été développées et comparées,
correspondant aux approches définies lors de la conception (Chapitre~4)~:

\begin{enumerate}
    \item \textbf{Modèle seul :} inférence directe du LLM sans contexte additionnel.
          Approche généraliste et rapide, servant de référence de comparaison.

    \item \textbf{Modèle avec ancrage juridique :} les passages de loi les plus
          pertinents sont récupérés dynamiquement et injectés dans la requête,
          ancrant la décision dans le droit français sur le harcèlement.

    \item \textbf{Règles heuristiques :} correspondance par mots-clés spécifiques
          au français, sans recours au LLM.
          Sert de mécanisme de repli rapide en cas d'indisponibilité du modèle.

    \item \textbf{Hybride :} appel au LLM avec repli automatique sur les heuristiques
          en cas d'échec ou d'ambiguïté de la réponse du modèle.
\end{enumerate}

L'implémentation est disponible en deux versions complémentaires~:
\begin{itemize}
    \item \textbf{Version Python :} adaptée aux expérimentations sur des corpus réduits,
          permettant de tester rapidement les quatre stratégies et de comparer
          leurs résultats avant mise en production.
    \item \textbf{Version Scala/Spark :} conçue pour le traitement en parallèle
          à grande échelle avec Apache Spark~3.5, distribuant l'analyse sur
          l'ensemble des messages soumis à la plateforme Netethic.
\end{itemize}

\subsection{Évaluation des performances}
\label{subsec:eval_classif}

L'évaluation comparative a été conduite en deux temps~:
une première évaluation globale sur le corpus Scala/Spark (1\,500~textes,
16~catégories), puis une évaluation ciblée sur un \textbf{échantillon stratifié}
de 150~textes (50~par classe~: Harcèlement, Santé mentale, Non-harcèlement)
afin d'obtenir des métriques équilibrées, non biaisées par la distribution
naturelle des classes.

\subsubsection*{Résultats globaux — pipeline Scala/Spark}

Le Tableau~\ref{tab:eval_classif} présente les métriques d'évaluation
des quatre stratégies sur le corpus de test de 1\,500~textes.

\vspace{6pt}
\begin{table}[H]
\centering
\caption{Évaluation comparative des stratégies — corpus Scala/Spark (1\,500 textes)}
\label{tab:eval_classif}
\begin{tabular}{|l|c|c|c|}
\hline
\textbf{Stratégie} & \textbf{Précision} & \textbf{Rappel} & \textbf{F1-score} \\
\hline
Modèle seul                   & 0{,}295 & 0{,}437 & 0{,}352 \\
\hline
Modèle avec ancrage juridique & 0{,}223 & 0{,}606 & 0{,}326 \\
\hline
Règles heuristiques           & 0{,}110 & 0{,}333 & 0{,}167 \\
\hline
Hybride                       & 0{,}631 & 0{,}607 & \textbf{0{,}619} \\
\hline
\end{tabular}
\end{table}

La stratégie hybride atteint le meilleur F1-score global (0{,}619),
soit une amélioration de~76\,\% par rapport au modèle seul (0{,}352).
La stratégie avec ancrage juridique améliore le rappel (0{,}606) au détriment
de la précision (0{,}223), traduisant une tendance à surclasser les textes
ambigus comme relevant du harcèlement.
Les règles heuristiques, en classant la quasi-totalité des textes dans la
catégorie Non-harcèlement, obtiennent un F1-score macro de 0{,}167 ~:
elles ne constituent qu'un mécanisme de repli et non une stratégie principale.

\subsubsection*{Évaluation ciblée — pipeline Python LLM (échantillon stratifié)}

Pour valider le comportement du pipeline Python sur les trois classes cibles
de la plateforme Netethic, une évaluation sur un \textbf{échantillon stratifié
de 150~textes} (50~par classe) a été conduite sur \textbf{cinq configurations}~:
la ligne de base heuristique, deux variantes du prompt RAG binaire (sans et avec
RAG juridique), et deux variantes du prompt Multilabel (sans et avec RAG).

\vspace{6pt}
\begin{table}[H]
\centering
\caption{Évaluation sur échantillon stratifié — pipeline Python LLM (150 textes, 50/classe)}
\label{tab:eval_python}
\small
\begin{tabular}{|l|l|c|c|c|}
\hline
\textbf{Stratégie} & \textbf{Classe} & \textbf{Précision} & \textbf{Rappel} & \textbf{F1-score} \\
\hline
\multirow{4}{*}{Règles heuristiques}
 & Harcèlement      & 0{,}00 & 0{,}00 & 0{,}00 \\
 & Santé mentale    & 0{,}00 & 0{,}00 & 0{,}00 \\
 & Non-harcèlement  & 0{,}33 & 1{,}00 & 0{,}50 \\
 & \textit{Macro}   & \textit{0{,}11} & \textit{0{,}33} & \textit{0{,}17} \\
\hline
\multirow{3}{*}{RAG-prompt (sans RAG)}
 & Harcèlement      & 0{,}69 & 0{,}70 & 0{,}69 \\
 & Santé mentale    & --- & --- & --- \\
 & Non-harcèlement  & 0{,}43 & 0{,}86 & 0{,}58 \\
 & \textit{Macro}   & \textit{0{,}37} & \textit{0{,}52} & \textit{0{,}42} \\
\hline
\multirow{3}{*}{RAG-prompt (avec RAG)}
 & Harcèlement      & 0{,}73 & 0{,}76 & 0{,}75 \\
 & Santé mentale    & --- & --- & --- \\
 & Non-harcèlement  & 0{,}45 & 0{,}88 & 0{,}59 \\
 & \textit{Macro}   & \textit{0{,}39} & \textit{0{,}55} & \textit{0{,}45} \\
\hline
\multirow{4}{*}{Multilabel LLM (sans RAG)}
 & Harcèlement      & 0{,}51 & 0{,}74 & 0{,}61 \\
 & Santé mentale    & \textbf{0{,}94} & 0{,}34 & 0{,}50 \\
 & Non-harcèlement  & 0{,}53 & 0{,}64 & 0{,}58 \\
 & \textit{Macro}   & \textit{0{,}66} & \textit{0{,}57} & \textit{0{,}56} \\
\hline
\multirow{4}{*}{\textbf{Multilabel LLM (avec RAG)}}
 & Harcèlement      & \textbf{0{,}68} & \textbf{0{,}76} & \textbf{0{,}72} \\
 & Santé mentale    & 0{,}71 & 0{,}40 & 0{,}51 \\
 & Non-harcèlement  & \textbf{0{,}55} & 0{,}72 & \textbf{0{,}62} \\
 & \textit{Macro}   & \textit{0{,}65} & \textit{0{,}63} & \textbf{\textit{0{,}62}} \\
\hline
\end{tabular}
\end{table}

\noindent\textit{Note~: le prompt RAG binaire prédit uniquement deux classes
(Harcèlement / Non-harcèlement)~; la Santé mentale est absente de sa taxonomie (---).}

La Figure~\ref{fig:eval_classif_comparison} illustre graphiquement
la comparaison entre les cinq stratégies sur les trois dimensions d'évaluation.

\begin{figure}[H]
\centering
\includegraphics[width=\textwidth]{figures/eval_classif_comparison.pdf}
\caption{Comparaison des cinq stratégies de classification~: métriques globales
         (F1 macro, accuracy), F1 par classe, et précision/rappel pour les
         classes Harcèlement et Santé mentale.
         Échantillon stratifié de 150~textes (50 par classe).}
\label{fig:eval_classif_comparison}
\end{figure}

La configuration \textbf{Multilabel LLM avec RAG} obtient le meilleur F1 macro
(\textbf{0{,}62}) et la meilleure exactitude (\textbf{63{,}0\,\%}),
confirmant l'apport combiné du prompt taxonomique et de l'ancrage RAG juridique.
Le prompt RAG binaire, bien qu'il atteigne un F1 de 0{,}75 sur la classe
Harcèlement avec RAG, ne peut structurellement pas détecter la Santé mentale,
ce qui pénalise son F1 macro (0{,}45).
L'analyse par classe pour la meilleure configuration révèle~:
\begin{itemize}
    \item \textbf{Harcèlement~:} rappel de \textbf{0{,}76} et F1 de 0{,}72 —
          le modèle détecte les trois quarts des cas avec une précision de 0{,}68,
          améliorée grâce aux exemples légaux récupérés par le RAG.
    \item \textbf{Santé mentale~:} précision de 0{,}71 et rappel de 0{,}40 —
          les signaux de détresse sont correctement identifiés lorsque détectés~;
          le rappel modéré reflète la difficulté à distinguer les expressions
          émotionnelles subtiles.
    \item \textbf{Non-harcèlement~:} F1 de \textbf{0{,}62} avec un rappel de 0{,}72 —
          le modèle reconnaît correctement la majorité des échanges neutres,
          limitant les faux positifs sur les conversations anodines.
\end{itemize}
Ces résultats valident l'architecture Multilabel LLM avec RAG comme stratégie
optimale pour la détection tripartite du harcèlement en ligne au sein de la
plateforme Netethic.

% ─────────────────────────────────────────────────────────────────────────────
\section{Conclusion}

Ce chapitre a présenté la phase de préparation des données et de mise en place
des modèles d'intelligence artificielle pour les quatre modules du système.

Pour les modules conversationnels, la qualité des réponses repose directement
sur la qualité des bases de connaissances~: les vingt-et-une pages publiques
de Netethic ont été transformées en 177~segments structurés selon une stratégie
parent-enfant pour Noa, et les documents du Support Assistant ont été organisés
en onze collections correspondant aux catégories fonctionnelles de la plateforme.
Dans les deux cas, un enrichissement manuel des fichiers Markdown — clarification
des sous-titres et complétion des informations manquantes — a permis d'améliorer
la pertinence des résultats de recherche vectorielle.

Pour le classificateur de harcèlement, le cycle d'apprentissage complet —
collecte de 977~textes, nettoyage en cinq étapes, augmentation par paraphrase
automatique, labeling indépendant et comparaison de quatre stratégies —
a abouti à un F1-score de 0{,}619 pour la stratégie hybride sur le corpus
Scala/Spark, contre 0{,}352 pour le modèle seul.
L'évaluation ciblée du pipeline Python LLM sur échantillon stratifié confirme
la supériorité de l'approche Multilabel LLM avec RAG (F1 macro \textbf{0{,}62})
sur toutes les autres configurations, dont le prompt RAG binaire (0{,}45) et la
ligne de base heuristique (0{,}17), avec des F1-scores de 0{,}72 sur Harcèlement,
0{,}51 sur Santé mentale et 0{,}62 sur Non-harcèlement.

Le chapitre suivant présente la réalisation concrète du système~:
interfaces utilisateur, résultats des tests automatisés, mesures de sécurité
et optimisations de performance.
